%================================
% Vocabulary & Language Focus — XB2 (Unit 4)
% 编译方式: xelatex main.tex (需运行两遍以生成目录)
%================================
\documentclass[../main.tex]{subfiles}

\begin{document}

\section{Unit 4}

\vp{(build) castles in the air}

\sense{v.}{（建）空中楼阁，幻想，空想}
\ex{Stop building castles in the air and start working.}
\ex{His plans for the future are just castles in the air.}

\sense{n.}{castle 城堡（拓展）}
\ex{The castle overlooks the valley.}
\col{a sandcastle 沙堡}

\begin{origin}
英语说 castles in the air，法语却说 \textit{châteaux en Espagne}（在西班牙建城堡）——中古时代西班牙是浪漫的远方，无法企及之地。castle 一词来自拉丁语 \textit{castellum}（小堡垒），是 \textit{castrum}（军营）的指小词——古罗马军团营地退役后演变成城镇与城堡。
\end{origin}

\vw{input}{/ˈɪnpʊt/}{n. \& vt.}

\sense{n.}{输入的信息}
\ex{The computer processes the input data.}

\sense{n.}{投入}
\col{input into/to sth}
\ex{The project needs more financial input.}
\col{input from sb 某人的意见（拓展）}
\ex{We need input from all departments.}

\sense{vt.}{(input, input) 输入（信息）}
\col{input data/a password}
\ex{Please input your password and press Enter.}

\begin{gram}[input 的不规则变化]
\textit{input} 的过去式与过去分词既可写 \textit{input}（同 put），也可写 \textit{inputted}——前者更常用。这类"由动词词组转化而来"的动词（in + put）常保留原动词的变化方式，但正规语法书多接受双写：\textit{The data was input/inputted.}
\end{gram}

\begin{origin}
in（进）+ put（放）——"放进去的东西"。20 世纪随计算机术语 output（输出）对举而产生：input 是"放进去"，output 是"放出来"。
\end{origin}

\vw{capacity}{/kəˈpæsəti/}{n.}

\sense{n.}{领悟（或理解、办事）能力}
\col{a capacity for sth}
\ex{She has a great capacity for learning languages.}

\sense{n.}{容量，容积}
\col{the capacity of sth}
\ex{The stadium has a seating capacity of 80,000.}
\col{at full capacity 满负荷}
\ex{The factory is running at full capacity.}

\sense{n.}{职位，职责}
\col{in one's capacity as sth 以……的身份}
\ex{He spoke in his capacity as chairman.}

\sense{n.}{生产量}
\col{production capacity 生产能力}
\ex{The company is expanding its production capacity.}

\begin{origin}
来自拉丁语 \textit{capacitas}，源于 \textit{capax}（能容纳的），词根 \textit{capere}（取、握）。同源词：\textbf{capable}（有能力的）、\textbf{capture}（捕获）、\textbf{captive}（俘虏）、\textbf{accept}（接受）。"能握住多少"——容量与能力，同一意象。
\end{origin}

\vw{massive}{/ˈmæsɪv/}{adj.}

\sense{adj.}{巨大的，大而重的}
\col{a massive building/rock}
\ex{The castle was built of massive stones.}

\sense{adj.}{非常严重的}
\col{a massive heart attack 严重的心脏病发作}
\col{massive losses/damage}
\ex{The company made massive losses last year.}

\col{massive changes 巨大的变化}

\begin{origin}
来自法语 \textit{massif}（实心的、笨重的），源于 \textit{masse}（块、堆），再追溯到拉丁语 \textit{massa}（面团、块）。同源词：\textbf{mass}（质量；群众）、\textbf{massif}（山块）。"一整块"自然"巨大沉重"。
\end{origin}

\vw{network}{/ˈnetwɜːk/}{n.}

\sense{n.}{网络，网状系统}
\col{a rail/road network 铁路网；公路网}
\ex{The country has an extensive rail network.}

\sense{n.}{关系网，人际网}
\col{a network of spies/friends}
\ex{He built up a network of contacts in the industry.}

\sense{n.}{（互联）网络}
\col{a social network 社交网络}
\ex{Social networks have changed how we communicate.}

\col{network (v.) 建立人脉（拓展）}
\ex{She's good at networking at conferences.}

\begin{origin}
net（网）+ work（结构）。16 世纪指网状织物（如蛛网般的蕾丝），20 世纪无线电与计算机时代引申为"网络"。汉语"上网"即 the Net / the network。
\end{origin}

\vw{coin}{/kɔɪn/}{vt.}

\sense{vt.}{创造（新词语）}
\col{coin a word/phrase/term}
\ex{The writer coined the term "cyberspace".}
\col{to coin a phrase 套用一句老话}
\ex{To coin a phrase, better late than never.}

\sense{n.}{硬币（拓展）}
\ex{He collects old coins.}

\begin{origin}
来自拉丁语 \textit{cuneus}（楔子）$\to$ 古法语 \textit{coing}（铸币的楔形模子）。硬币由楔形模具"铸"出，故"铸币"与"造词"共用一词——coin a word 即"铸造一个新词"。同源词：\textbf{cuneiform}（楔形文字）。
\end{origin}

\vw{breakthrough}{/ˈbreɪkθruː/}{n. \& adj.}

\sense{n.}{重大进展，突破}
\col{a major breakthrough}
\ex{The discovery is a major breakthrough in cancer research.}
\col{make/achieve a breakthrough}
\ex{Scientists have made a breakthrough in battery technology.}

\sense{adj.}{突破性的}
\col{a breakthrough deal/agreement}
\ex{The talks produced a breakthrough agreement.}

\begin{origin}
break（突破）+ through（穿过）——"破围而出"。本是二战军事术语"突围"，战后转用于科技、谈判领域。
\end{origin}

\vw{finance}{/ˈfaɪnæns; faɪˈnæns/}{n. \& vt.}

\sense{n.}{财政，金融}
\col{the finance industry 金融业}
\ex{He works in finance in London.}

\sense{n.}{资金}
\ex{The project needs external finance.}

\sense{n.}{财力，财务管理}
\col{personal finances 个人财务状况}
\ex{He manages the household finances.}

\sense{vt.}{提供资金}
\ex{The new airport was financed by the government.}

\col{financial (adj.) 财政的}
\ex{financial support}

\begin{origin}
来自古法语 \textit{finance}（结清、付款），源于 \textit{finer}（结束、结清），词根 \textit{fin}（结束）。同源词：\textbf{finish}（完成）、\textbf{final}、\textbf{fine}（罚金）。"把钱结清"就是金融——finance 与 finish 同源。
\end{origin}

\vp{to name but a few}

\sense{adv.}{略举几例}
\ex{The museum has works by Picasso, Matisse and Dalí, to name but a few.}
\col{同义：among others}

\begin{origin}
name（举名）+ but（仅仅）——"只举几个名字"，表示举例远不止于此。
\end{origin}

\vw{autonomous}{/ɔːˈtɒnəməs/}{adj.}

\sense{adj.}{自主的，有自主权的；自治的}
\col{an autonomous region 自治区}
\ex{The area enjoys autonomous status.}
\col{autonomous vehicles 自动驾驶汽车}
\ex{Autonomous cars may soon be on our roads.}

\col{autonomy (n.) 自治；自主权}

\begin{origin}
来自希腊语 \textit{autonomos}（自治的）：\textit{autos}（自己）+ \textit{nomos}（法律）。"自己给自己立法"就是自治。同源词：\textbf{economy}（oikos 家 + nomos 管理 $\to$ 治家 $\to$ 经济）——economy 与 autonomous 共享 nomos 一词根。
\end{origin}

\vw{domestic}{/dəˈmestɪk/}{adj.}

\sense{adj.}{家庭的，家用的}
\col{domestic life 家庭生活}
\ex{The new job left him little time for domestic life.}
\col{domestic violence 家庭暴力}
\col{domestic animals 家畜}

\sense{adj.}{本国的，国内的}
\col{domestic flights 国内航班}
\ex{The airline operates both domestic and international flights.}
\col{the domestic market 国内市场}

\begin{origin}
来自拉丁语 \textit{domesticus}（家的），词根 \textit{domus}（家、房屋）。同源词：\textbf{dome}（圆顶——屋宇之顶）、\textbf{domicile}（住所）、\textbf{major-domo}（管家）。“家的事”扩展为“国内的事”。
\end{origin}

\vw{threaten}{/ˈθretn/}{vt.}

\sense{vt.}{危及，对……构成威胁}
\ex{The factory threatens the local wildlife.}
\ex{The flood threatened the whole village.}

\sense{vt.}{威胁，恐吓}
\col{threaten sb with sth}
\ex{He threatened the driver with a knife.}
\col{threaten to do sth}
\ex{The workers threatened to go on strike.}

\col{threatening (adj.) 威胁性的}
\ex{a threatening letter}

\begin{origin}
threat（威胁）源自古英语 \textit{þrēat}（人群、压迫、威胁），与"拥挤压迫"的意象有关。同源词：\textbf{threat}、\textbf{threaten}。日耳曼语词根。
\end{origin}

\vw{commit}{/kəˈmɪt/}{vt. \& vi.}

\sense{vt.}{犯（罪）}
\col{commit a crime/murder}
\ex{He committed a serious crime and was jailed.}
\col{commit suicide 自杀}

\sense{vt.}{承诺，保证}
\col{commit to sth/doing sth}
\ex{The government committed to reducing emissions.}
\col{commit oneself to sth 承诺；献身}
\ex{She committed herself to the cause of education.}

\sense{vt.}{花（钱或时间）}
\col{commit money/time to sth}
\ex{The company committed \$2 million to the project.}

\sense{vi.}{全心全意投入，忠于}
\col{commit to sth}
\ex{We are fully committed to this partnership.}

\begin{origin}
来自拉丁语 \textit{committere}（托付、送交）：\textit{com-}（一起）+ \textit{mittere}（送出）。同源词：\textbf{mission}（使命）、\textbf{transmit}（传送）、\textbf{submit}（提交）、\textbf{admit}（承认）、\textbf{commission}（委托）。"把自己送出去"是投入，把"罪行送出去"是犯罪——同一动作，正邪两途。（committed、commitment 见 Unit 1。）
\end{origin}

\vw{crime}{/kraɪm/}{n.}

\sense{n.}{罪，罪行；犯罪活动}
\col{commit a crime 犯罪}
\ex{The police caught the man who committed the crime.}
\col{the crime rate 犯罪率}
\ex{The crime rate has fallen this year.}
\col{organized crime 有组织犯罪}
\col{a crime against humanity 反人类罪}

\col{criminal (adj. \& n.) 犯罪的；罪犯}

\begin{origin}
来自拉丁语 \textit{crimen}（指控、罪行），源于 \textit{cernere}（筛选、审判）——与 Unit 1 的 \textbf{discrimination}（discernere 辨别）同根。"经过审判断定的行为"就是罪行。
\end{origin}

\vw{murder}{/ˈmɜːdə(r)/}{vt. \& n.}

\sense{vt. \& n.}{谋杀，凶杀}
\col{murder (vt.) 谋杀}
\ex{He was accused of murdering his wife.}
\col{murder (n.) 谋杀}
\ex{The police are investigating the murder.}
\col{attempted murder 谋杀未遂}
\col{a murderer 凶手}

\begin{origin}
源自古英语 \textit{morthor}（谋杀），日耳曼语，词根与"死亡"有关——拉丁语 \textit{mors}（死）同源。同源词：\textbf{mortal}（凡人的；致命的）、\textbf{mortgage}（mort 死 + gage 抵押 $\to$ "死抵"：房产抵押直到还清）。murder 与 mortgage 共享"死"的词根。
\end{origin}

\vw{urge}{/ɜːdʒ/}{vt. \& n.}

\sense{vt.}{敦促，催促}
\col{urge sb to do sth}
\ex{The UN urged both sides to negotiate.}
\col{urge that ...（should）do（正式）}
\ex{They urged that the plan (should) be revised.}

\sense{vt.}{大力推荐，竭力主张}
\col{urge sth}
\ex{Doctors urge regular exercise.}

\sense{vt.}{驱赶，鞭策}
\col{urge sb on 为某人鼓劲}
\ex{The crowd urged the runner on.}

\sense{n.}{冲动，强烈的欲望}
\col{an urge to do sth}
\ex{He felt a sudden urge to laugh.}

\col{urgent (adj.) 紧急的}

\begin{origin}
来自拉丁语 \textit{urgere}（推、逼迫）。同源词：\textbf{urgent}（紧急的——逼迫人的事）。“推着人往前走”就是敦促；心中"被推着走"的感觉就是冲动。
\end{origin}

\vw{pose}{/pəʊz/}{vt. \& vi. \& n.}

\sense{vt.}{造成，产生}
\col{pose a threat/risk to sth}
\ex{The flood posed a serious threat to the village.}
\col{pose a question 提问}
\ex{The report poses important questions for policymakers.}

\sense{vi.}{摆好姿势}
\col{pose for a photo}
\ex{They posed for a group photo.}
\col{pose as sb 冒充}
\ex{He posed as a police officer.}

\sense{n.}{姿势}
\col{strike a pose 摆造型}

\begin{origin}
来自拉丁语 \textit{pausare}（停止、休息）$\to$ 古法语 \textit{poser}（放置）——poser 兼"放置"与"暂停"两义，由此分出英语 pose 与 pause 两个词（同为拉丁 pausare 的双重词）。"摆出姿态"是"放置自己"，"摆出问题"是"放置疑问"。
\end{origin}

\vw{apparently}{/əˈpærəntli/}{adv.}

\sense{adv.}{显然，可见}
\ex{Apparently, the meeting has been cancelled.}
\ex{The house was apparently empty.}
\col{apparent (adj.) 明显的}
\ex{It soon became apparent that we were lost.}

\begin{origin}
来自拉丁语 \textit{apparens}（显现的），源于 \textit{apparere}（显现）：\textit{ad-}（向）+ \textit{parere}（出现）。同源词：\textbf{appear}（出现）、\textbf{appearance}（外表）、\textbf{disappear}（消失）。“显现出来的”自然"显然"。
\end{origin}

\vw{warning}{/ˈwɔːnɪŋ/}{n.}

\sense{n.}{警告，先兆；警戒}
\col{give/send a warning}
\ex{The doctor gave me a warning about my diet.}
\col{a warning sign 警示标志；先兆}
\ex{The sudden drop in sales is a warning sign.}
\col{without warning 毫无预警}
\ex{The storm hit the coast without warning.}

\begin{origin}
warn 源自古英语 \textit{warnian}（提防、警示），日耳曼语。同源词：\textbf{wary}（谨慎的）、\textbf{aware}（察觉的）。"让人预先提防"——warn 与 aware 同根。
\end{origin}

\vw{demand}{/dɪˈmɑːnd/}{vt. \& n.}

\sense{vt.}{强烈要求}
\col{demand sth}
\ex{The workers demanded higher wages.}
\col{demand to do sth}
\ex{She demanded to see the manager.}

\sense{vt.}{需要}
\ex{This task demands great patience.}

\sense{n.}{要求，所需之物}
\col{meet/satisfy a demand 满足需求}
\ex{The company struggled to meet demand.}

\sense{n.}{需求，需要}
\col{in demand 抢手，需求大}
\ex{Good engineers are in high demand.}
\col{on demand 按需}
\ex{video on demand}

\begin{origin}
来自拉丁语 \textit{demandare}（委托、吩咐）：\textit{de-}（加强）+ \textit{mandare}（吩咐）。同源词：\textbf{command}（命令）、\textbf{mandate}（授权）、\textbf{recommend}（推荐）。“吩咐、要求”是 demand 的本义。
\end{origin}

\vw{victim}{/ˈvɪktɪm/}{n.}

\sense{n.}{受害者，牺牲品；受骗者}
\col{the victim of a crime/robbery}
\ex{She was the victim of a robbery.}
\col{fall victim to 受伤，受损，被害}
\ex{Many small shops fell victim to the recession.}
\col{a victim of fraud 诈骗案的受骗者}

\sense{n.}{祭品}
\col{a sacrificial victim 祭品}

\begin{origin}
来自拉丁语 \textit{victima}（祭品）。古罗马祭祀中献给神的牺牲，即是"受害者"——victim 本义是"献祭的羔羊"。同源词：\textbf{victimize}（使受害）。
\end{origin}

\vw{recognition}{/ˌrekəɡˈnɪʃn/}{n.}

\sense{n.}{识别，认出}
\col{beyond recognition 认不出来}
\ex{The town has changed beyond recognition.}

\sense{n.}{承认，认可}
\col{recognition of sth}
\ex{The agreement marks official recognition of their rights.}

\sense{n.}{赞誉，赏识}
\col{in recognition of sth 表彰，酬谢}
\ex{He received a medal in recognition of his bravery.}
\col{gain international recognition 获得国际认可}

\col{recognize (v.) 认出；承认}

\begin{origin}
来自拉丁语 \textit{recognoscere}（重新知道）：\textit{re-}（再）+ \textit{cognoscere}（认识），词根可追溯到原始印欧语 *\textit{gno-}（知道）——与英语 \textbf{know} 同源。recognition 与 know，跨越三千年仍是同根词。
\end{origin}

\vw{alternative}{/ɔːlˈtɜːnətɪv/}{n. \& adj.}

\sense{n.}{可供选择的事物}
\col{an alternative to sth}
\ex{There is no alternative to hard work.}

\sense{adj.}{可供替代的}
\col{alternative sources of energy 替代能源}
\ex{We need to develop alternative sources of energy.}

\sense{adj.}{非传统的}
\col{alternative music 另类音乐}
\ex{He listens to alternative rock.}

\begin{origin}
来自拉丁语 \textit{alternare}（交替），源于 \textit{alternus}（一个接一个的），词根 \textit{alter}（另一个）。同源词：\textbf{alter}（改变——变成另一个）、\textbf{alternate}（交替的）。"另一个选项"就是 alternative。
\end{origin}

\vp{break into}

\sense{v.}{强行闯入，撬开}
\ex{Burglars broke into the house last night.}
\col{break into a computer system 侵入系统}

\sense{v.}{突然开始（笑、唱等）}
\col{break into laughter/song/applause}
\ex{The audience broke into applause.}

\sense{v.}{顺利打入}
\col{break into a market}
\ex{The brand is trying to break into the Asian market.}

\begin{origin}
break（破）+ into（进入）——"破门而入"。三个义项共享"突然进入某状态"的意象：闯入房间、突然笑出声、打入市场。
\end{origin}

\vw{fail-safe}{/ˈfeɪl seɪf/}{adj.}

\sense{adj.}{有自动保险装置的，具有自动防止故障性能的}
\col{a fail-safe system/device}
\ex{The nuclear plant is equipped with a fail-safe system.}

\begin{origin}
fail（失灵）+ safe（安全）——"即使失灵也安全"。20 世纪工程术语：系统故障时自动进入安全状态，如电梯的制动器、飞机的自动驾驶。
\end{origin}

\vw{distribution}{/ˌdɪstrɪˈbjuːʃn/}{n.}

\sense{n.}{分发，分送}
\col{the distribution of sth}
\ex{The Red Cross organized the distribution of food and water.}

\sense{n.}{分配，分布}
\col{the distribution of wealth 财富分配}
\ex{The distribution of income is highly uneven.}
\col{the distribution of population 人口分布}

\col{distribution network 分销网络}
\col{distribute (v.) 分发；分布}

\begin{origin}
来自拉丁语 \textit{distribuere}（分给）：\textit{dis-}（分开）+ \textit{tribuere}（给、分配）——古罗马按"部落"（tribus）划分区域分配物资。同源词：\textbf{tribute}（贡品）、\textbf{contribute}（贡献：一起给）、\textbf{attribute}（归于）。
\end{origin}

\vw{predict}{/prɪˈdɪkt/}{vt.}

\sense{vt.}{预言，预告}
\col{predict sth}
\ex{Experts predict a rise in prices.}
\col{predict that ...}
\ex{She predicted that the company would go bankrupt.}
\col{be predicted to do sth}
\ex{The match is predicted to be very close.}

\col{prediction (n.) 预言}

\begin{origin}
来自拉丁语 \textit{praedicere}（预先说）：\textit{prae-}（预先）+ \textit{dicere}（说）。同源词：\textbf{dictate}（口述）、\textbf{dictionary}、\textbf{contradict}（反驳：对着说）。"提前说"就是预言。
\end{origin}

\vw{decrease}{/dɪˈkriːs/ \ vt. \& vi. \ \ /ˈdiːkriːs/ \ n.}{vt. \& vi. \& n.}

\sense{vt. \& vi. /dɪˈkriːs/}{减少，降低}
\col{decrease by/to sth}
\ex{The number of visitors decreased by 10\% last year.}
\ex{Temperature decreased to below zero at night.}

\sense{n. /ˈdiːkriːs/}{减少，降低}
\col{a decrease in sth}
\ex{We saw a decrease in sales this quarter.}

\begin{gram}[名动异读：重音转移]
\textit{decrease} 名词重音在前（/\textbf{ˈdiː}kriːs/），动词重音在后（/dɪ\textbf{ˈkriːs}/）——与 \textbf{contract}、\textbf{record}、\textbf{increase} 同类（见 Unit 3 contract 条）。记忆：重音在前多为名词。
\end{gram}

\begin{origin}
来自拉丁语 \textit{decrescere}（减少）：\textit{de-}（向下）+ \textit{crescere}（生长）——"反向生长"。同源词：\textbf{increase}（in- 向上 + crescere $\to$ 增长）、\textbf{crescent}（新月：正在生长的月亮）、\textbf{crew}（crescere 的派生：长成的团队）。
\end{origin}

\vw{deny}{/dɪˈnaɪ/}{vt.}

\sense{vt.}{否认，否定}
\col{deny (doing) sth}
\ex{He denied stealing the money.}
\col{deny that ...}
\ex{She denied that she had said it.}

\sense{vt.}{拒绝承认，拒绝接受}
\col{deny sb sth 拒绝给予}
\ex{The boy was denied access to the library.}
\col{There is no denying ... 不可否认}
\ex{There is no denying that he is talented.}

\begin{lex}{deny / refuse / reject}
\textbf{deny} 表"否认事实"或"拒绝给予"（\textit{deny a rumour} 辟谣；\textit{deny sb permission}）；\textbf{refuse} 表"拒绝去做"（\textit{refuse to answer}）；\textbf{reject} 表"拒不接受"（\textit{reject an offer}）。句型差异：\textit{deny doing}；\textit{refuse to do}；\textit{reject sth}。
\end{lex}

\begin{origin}
来自拉丁语 \textit{denegare}（拒绝）：\textit{de-}（完全）+ \textit{negare}（说不）。同源词：\textbf{negative}（否定的）、\textbf{negate}（否定）。"说不是"就是否认。
\end{origin}

\vw{quantity}{/ˈkwɒntəti/}{n.}

\sense{n.}{数量}
\col{a quantity of sth 一些，若干}
\ex{Add a small quantity of salt.}

\sense{n.}{大量}
\col{in large quantities 大量，大批}
\ex{The factory produces goods in large quantities.}
\col{in quantity 大量地}
\ex{The wine is sold in quantity to restaurants.}

\begin{gram}[quantity 的主谓一致]
\textit{a quantity of} 后接不可数名词，谓语用单数：\textit{A large quantity of water was wasted.}；\textit{quantities of} 后接可数复数或不可数名词，谓语用复数：\textit{Quantities of books are sold online.} 考试常考这一组。
\end{gram}

\begin{origin}
来自拉丁语 \textit{quantitas}（多少），源于 \textit{quantus}（多大、多少）。同源词：\textbf{quantum}（量子——一份一份的量）、\textbf{quantity}。"有多少"就是数量。
\end{origin}

\vw{airline}{/ˈeəlaɪn/}{n.}

\sense{n.}{航空公司}
\ex{The airline operates flights to 50 countries.}
\col{airline tickets 机票}

\begin{origin}
air（空中）+ line（航线）——"空中的航线"。line 本义"线"，引申为"航路"（\textit{shipping line} 航运公司亦此构词）。20 世纪初航空业诞生后，"航空公司"沿用这一命名模式。
\end{origin}

\vw{threat}{/θret/}{n.}

\sense{n.}{构成威胁的人（或事物）}
\col{a threat to sth}
\ex{Pollution is a serious threat to wildlife.}

\sense{n.}{威胁，恐吓}
\col{make threats against sb}
\ex{He made threats against the witness.}

\sense{n.}{凶兆}
\col{under threat 受到威胁}
\ex{The forest is under threat from logging.}
\col{a threat of rain 要下雨的迹象}

\begin{origin}
源自古英语 \textit{þrēat}（人群、压迫、威胁）。同源词：\textbf{threaten}（威胁）。"压迫性的力量"——威胁与恐吓同出一源。
\end{origin}

\vw{spy}{/spaɪ/}{vi. \& vt. \& n.}

\sense{vi.}{从事间谍活动}
\col{spy on sb 暗中监视，窥探}
\ex{She suspected her neighbour was spying on her.}

\sense{vt.}{突然看见}
\ex{I spied a familiar face in the crowd.}

\sense{n.}{间谍}
\ex{He worked as a spy during the war.}

\col{espionage (n.) 谍报活动}

\begin{origin}
来自古法语 \textit{espier}（窥视），源于日耳曼语 *\textit{speh-}（看）——与拉丁语 \textit{specere}（看）远亲。同源词：\textbf{espionage}（谍报）、\textbf{spectator}（观众）、\textbf{spyglass}（望远镜）。"暗中看的人"就是间谍。
\end{origin}

\vw{monitor}{/ˈmɒnɪtə(r)/}{vt. \& n.}

\sense{vt.}{监视，跟踪调查；监听}
\col{monitor air quality/progress}
\ex{The equipment monitors air quality around the clock.}
\col{monitor the situation 密切关注局势}
\ex{We are monitoring the situation closely.}

\sense{n.}{显示屏；监测器}
\col{a computer monitor 显示器}
\ex{He stared at the monitor all day.}
\col{a heart monitor 心脏监测器}

\sense{n.}{班长}
\col{the class monitor}
\ex{The class monitor collected the homework.}

\begin{origin}
来自拉丁语 \textit{monere}（提醒、警告）$\to$ \textit{monitor}（提醒者）。同源词：\textbf{monument}（纪念碑——提醒后人者）、\textbf{admonish}（告诫）、\textbf{monition}（警告）。显示器、监测器、班长——都是"看着、提醒着"的角色。
\end{origin}

\vw{highlight}{/ˈhaɪlaɪt/}{n. \& vt.}

\sense{n.}{最好（或最精彩）的部分}
\col{a highlight of sth}
\ex{The visit to the museum was the highlight of our trip.}
\col{the highlights 精彩集锦}
\ex{Watch the match highlights online.}

\sense{vt.}{突出，强调}
\ex{The report highlights the need for reform.}

\sense{vt.}{使醒目}
\col{highlight the text 高亮文本}
\ex{Highlight the key words in yellow.}

\begin{origin}
high（高）+ light（光）——"最亮的光"。19 世纪绘画术语"高光"（画面最亮处），后引申为"最精彩的部分"与"使突出"。
\end{origin}

\vw{capsule}{/ˈkæpsjuːl/}{n.}

\sense{n.}{小塑料容器}
\col{a time capsule 时间胶囊}
\ex{The school buried a time capsule.}

\sense{n.}{（装药的）胶囊}
\ex{The medicine comes in capsules.}

\sense{n.}{太空舱}
\col{a space capsule 太空舱}
\ex{The capsule returned safely to Earth.}

\begin{origin}
来自拉丁语 \textit{capsula}（小盒子），是 \textit{capsa}（盒子）的指小词，词根 \textit{capere}（握住）。同源词：\textbf{case}（箱子）、\textbf{capture}（捕获）、\textbf{capacity}。"能握住东西的小盒子"——药丸、太空舱，皆"盒"。
\end{origin}

\vw{bury}{/ˈberi/}{vt.}

\sense{vt.}{埋藏；埋葬，安葬}
\col{bury sth 埋藏}
\ex{The dog buried a bone in the garden.}
\ex{He was buried in the village churchyard.}

\sense{vt.}{覆盖；遮盖}
\col{be buried in sth 被……覆盖}
\ex{The house was buried in snow.}
\col{bury one's face in one's hands 掩面}
\ex{She buried her face in her hands and wept.}

\sense{vt.}{使陷入}
\col{be buried in work/thought 埋头工作；陷入沉思}
\ex{He was buried in his paperwork.}

\col{bury the hatchet 言归于好}

\begin{origin}
源自古英语 \textit{byrgan}（埋葬、藏），日耳曼语，与"保护、遮蔽"的意象有关。同源词：\textbf{burrow}（地洞——藏身之处）、\textbf{borough}（堡：庇护之所）。
\end{origin}

\vw{fountain}{/ˈfaʊntən/}{n.}

\sense{n.}{喷射，涌流；喷水池}
\col{a drinking fountain 饮水器}
\ex{The children played in the fountain.}

\sense{n.}{源泉，丰富来源}
\col{a fountain of sth}
\ex{He is a fountain of ideas.}
\col{the fountain of youth 青春之泉}

\col{fountain pen 自来水笔}
\ex{She wrote the letter with a fountain pen.}

\begin{origin}
来自拉丁语 \textit{fons}（泉；宾格 \textit{fontem}）$\to$ 古法语 \textit{fontaine}。同源词：\textbf{font}（字体：一套字模"泉涌"般源源不绝）。"泉"引申为"喷泉"与"源泉"。
\end{origin}

\vw{contemporary}{/kənˈtemprəri/}{adj. \& n.}

\sense{adj.}{当代的，现代的}
\col{contemporary art/music}
\ex{The museum specializes in contemporary art.}

\sense{adj.}{属同时期的，同一时代的}
\col{be contemporary with sb/sth}
\ex{Shakespeare was contemporary with Marlowe.}

\sense{n.}{同代人，同辈人}
\col{one's contemporaries}
\ex{He was admired by his contemporaries.}

\begin{origin}
来自拉丁语 \textit{contemporarius}：\textit{con-}（同）+ \textit{tempus}（时间）。同源词：\textbf{temporary}（暂时的）、\textbf{tempo}（节奏）、\textbf{tense}（时态）。"同一时间的人"——同代人，乃至"我们这个时代的人"。
\end{origin}

\vw{seed}{/siːd/}{n.}

\sense{n.}{种子}
\col{plant/sow seeds 播种}
\ex{The farmer sowed the seeds in spring.}
\col{grow from seed 从种子长起}

\sense{n.}{起源，起因}
\col{the seed of sth}
\ex{The trip planted the seeds of her interest in archaeology.}

\begin{origin}
源自古英语 \textit{sæd}（种子），日耳曼语，词根可追溯到原始印欧语 *\textit{se-}（播种）。同源词：\textbf{sow}（播种）、\textbf{season}（季节——原义"播种之时"，拉丁语 \textit{satio} 播种 $\to$ \textit{seison} $\to$ season）。seed 与 season 同源。
\end{origin}

\vw{muscular}{/ˈmʌskjələ(r)/}{adj.}

\sense{adj.}{肌肉的；强壮的}
\col{a muscular body/build}
\ex{The athlete has a muscular build.}
\col{muscular strength 肌肉力量}
\ex{The boxer relies on muscular strength.}

\begin{origin}
muscle 来自拉丁语 \textit{musculus}——"小老鼠"（\textit{mus} 老鼠的指小词）。罗马人觉得肱二头肌收缩时像一只小老鼠在皮下跑动，遂名 muscle。同源词：\textbf{mouse}（老鼠）、\textbf{muscular}（肌肉发达的）。肌肉与老鼠，同根相称三千年。
\end{origin}

\vw{conclude}{/kənˈkluːd/}{vt. \& vi.}

\sense{vt.}{断定，推断出}
\col{conclude that ...}
\ex{The report concluded that the plan was unworkable.}

\sense{vt.}{达成，订立}
\col{conclude an agreement/treaty}
\ex{The two sides concluded a peace treaty.}

\sense{vt. \& vi.}{（使）结束，终止}
\ex{The meeting concluded with a vote.}

\col{conclusion (n.) 结论（见 Unit 1）}

\begin{origin}
来自拉丁语 \textit{concludere}（关闭）：\textit{con-}（一起）+ \textit{claudere}（关闭）。同源词：\textbf{conclusion}、\textbf{include}、\textbf{exclude}、\textbf{close}。"把话说拢"就是下结论、收尾。
\end{origin}

\vw{justified}{/ˈdʒʌstɪfaɪd/}{adj.}

\sense{adj.}{合乎情理的，事出有因的；有正当理由的}
\col{be justified in doing sth 做……是合理的}
\ex{She was fully justified in complaining.}
\col{a justified decision/action}
\ex{The court upheld the decision as justified.}

\col{justify (v.) 证明……正当}
\ex{How can you justify such expense?}

\begin{origin}
来自拉丁语 \textit{iustificare}（使公正）：\textit{iustus}（公正的）+ \textit{facere}（做）。同源词：\textbf{just}（公正的）、\textbf{justice}（正义）、\textbf{adjust}（调整）。“做在理上”就是有正当理由。
\end{origin}

\vw{superiority}{/suːˌpɪəriˈɒrəti/}{n.}

\sense{n.}{优越感}
\col{a sense/air of superiority}
\ex{He spoke with an air of superiority.}

\sense{n.}{优越性，优势}
\col{superiority over sth}
\ex{The team proved its superiority over all rivals.}

\col{superior (adj.) 更好的；上级的}

\begin{origin}
来自拉丁语 \textit{superior}（更高的）——\textit{superus}（在上面的）的比较级，最高级是 \textit{supremus}。同源词：\textbf{super}（超级）、\textbf{supreme}（至高的）。“高别人一头”就是优越。
\end{origin}

\vw{virtual}{/ˈvɜːtʃuəl/}{adj.}

\sense{adj.}{虚拟的，模拟的}
\col{virtual reality 虚拟现实}
\ex{The game uses virtual reality technology.}
\col{a virtual tour 虚拟游览}

\sense{adj.}{很接近的，事实上的}
\col{a virtual certainty 几乎确定的事}
\ex{It's a virtual certainty that prices will rise.}
\col{virtually (adv.) 几乎，事实上}
\ex{The work is virtually finished.}

\begin{origin}
来自拉丁语 \textit{virtus}（美德、力量）$\to$ 中世纪经院哲学的 \textit{virtualis}——"虽有实效、但无实体"（如 virtus 的作用）。20 世纪计算机术语由此引申出"虚拟"：virtual reality 即"无实体而有实效的现实"。同源词：\textbf{virtue}（美德）。virtual 与 virtue 同源。
\end{origin}

\vp{at the forefront of}

\sense{prep.}{处于最前列，进入重要地位}
\ex{The university is at the forefront of AI research.}
\ex{The company is at the forefront of green technology.}

\begin{origin}
fore（前）+ front（正面）——"最前面的一排"。fore- 是英语古老的前缀（\textbf{forecast} 预报、\textbf{forehead} 额头），front 来自拉丁语 \textit{frons}（前额）。
\end{origin}

\vw{definition}{/ˌdefɪˈnɪʃn/}{n.}

\sense{n.}{定义；释义，解释}
\col{the definition of sth}
\ex{What is the dictionary definition of this word?}
\col{by definition 按定义；理所当然}
\ex{War is by definition destructive.}

\col{define (v.) 下定义}

\begin{origin}
来自拉丁语 \textit{definitio}（界限、定义），源于 \textit{definire}（划定界限）：\textit{de-}（完全）+ \textit{finire}（限定），词根 \textit{finis}（边界）。同源词：\textbf{finish}（完成）、\textbf{finite}（有限的）、\textbf{final}。“划定边界”就是下定义——definition 与 finish 同源。
\end{origin}

\vw{sensor}{/ˈsensə(r)/}{n.}

\sense{n.}{传感器，敏感元件}
\col{a light/motion sensor 光敏传感器；运动传感器}
\ex{The sensors detect any movement in the room.}

\begin{origin}
来自拉丁语 \textit{sentire}（感觉）$\to$ \textit{sensor}（感觉器）。同源词：\textbf{sense}（感觉）、\textbf{sensory}（感官的）、\textbf{sensitive}（敏感的）。"会感觉的元件"就是传感器。
\end{origin}

\vw{practical}{/ˈpræktɪkl/}{adj.}

\sense{adj.}{实际的，真实的}
\col{practical experience 实践经验}
\ex{The course offers plenty of practical training.}
\col{for practical purposes 实际上}

\sense{adj.}{切实可行的；有用的}
\col{a practical solution/suggestion}
\ex{We need a practical solution to the problem.}

\sense{adj.}{实事求是的}
\ex{She's a very practical person.}

\begin{origin}
来自希腊语 \textit{praktikos}（行动的），源于 \textit{prassein}（做）。同源词：\textbf{practice}（实践）、\textbf{pragmatic}（务实的——希腊语 pragma 事）。practical 与 pragmatic 同根："动手做的、管用的"。
\end{origin}

\vp{a host of}

\sense{det.}{许多，大量}
\ex{The book raises a host of questions.}
\ex{The plan has a host of benefits.}

\begin{origin}
host 有两个词源：作"主人"讲源自拉丁语 \textit{hospes}（客人、主人）；作"一大群"讲源自拉丁语 \textit{hostis}（敌人、军队）——"敌军"引申为"大队人马、许多"。同源词：\textbf{hostile}（敌对的）、\textbf{hostage}（人质）。a host of 用的是"军队"这支词源。
\end{origin}

\vw{academic}{/ˌækəˈdemɪk/}{adj.}

\sense{adj.}{教学的，学术的}
\col{academic research 学术研究}
\ex{He devoted himself to academic research.}
\col{the academic year 学年}

\sense{adj.}{学习良好的}
\ex{She has always been an academic child.}

\sense{adj.}{纯理论的}
\col{purely academic 纯属理论上的}
\ex{The debate is purely academic; nothing will change.}

\begin{origin}
来自希腊语 \textit{Akadēmeia}——柏拉图在雅典城外讲学的那片园林，得名于英雄 \textit{Akademos}。同源词：\textbf{academy}（学园；学院）。"学园里的"就是学术的。
\end{origin}

\vw{immersive}{/ɪˈmɜːsɪv/}{adj.}

\sense{adj.}{（计算机系统或图像）沉浸式虚拟现实的}
\col{an immersive experience 沉浸式体验}
\ex{The game offers a fully immersive experience.}
\col{immersive technology 沉浸式技术}

\begin{origin}
immerse 来自拉丁语 \textit{immergere}（浸入）：\textit{in-}（入）+ \textit{mergere}（浸、沉）。同源词：\textbf{merge}（合并——浸入彼此）、\textbf{submerge}（淹没）、\textbf{emerge}（浮现——从水中出来）。"把人浸入虚拟世界"就是沉浸式——immersive 与 emerge 同根。
\end{origin}

\vw{engage}{/ɪnˈɡeɪdʒ/}{vt. \& vi.}

\sense{vt.}{吸引住（注意力、兴趣）}
\col{engage sb's attention/interest}
\ex{The story engaged the children's interest immediately.}

\sense{vt.}{雇用，聘用}
\col{engage sb as sth}
\ex{They engaged a lawyer to handle the case.}

\sense{vt. \& vi.}{与……建立密切关系}
\col{engage in sth 从事，参与}
\ex{He engaged in politics after graduation.}
\col{engage with sb/sth 与……互动}
\ex{The teacher knows how to engage with her students.}

\col{be engaged to sb 订婚}
\ex{She is engaged to a doctor.}
\col{engagement (n.) 订婚；参与}

\begin{origin}
来自古法语 \textit{engagier}（抵押、许下）：\textit{en-}（入）+ \textit{gage}（抵押物）——与 \textbf{mortgage}（抵押贷款）同源。"把自己抵押出去"：订婚是许下终身，投入是全身心进入，聘用是雇人承诺。engage 与 mortgage 同根。
\end{origin}

\vw{barrier}{/ˈbæriə(r)/}{n.}

\sense{n.}{障碍，阻力}
\col{a barrier to sth}
\ex{Language is a barrier to communication.}
\col{a trade barrier 贸易壁垒}

\sense{n.}{屏障，障碍物}
\col{a crash barrier 防撞栏}
\col{the sound barrier 音障}

\sense{n.}{分界线，隔阂}
\col{break down barriers 消除隔阂}
\ex{Music can break down barriers between cultures.}

\begin{origin}
来自古法语 \textit{barriere}（门栏），源于 \textit{barre}（横木）。同源词：\textbf{bar}（横木；酒吧）、\textbf{barricade}（街垒）。“横在路上的木条”就是障碍。
\end{origin}

\vp{boil down to}

\sense{v.}{归结为，基本问题是}
\ex{It all boils down to money in the end.}
\ex{The problem boils down to a lack of trust.}

\begin{origin}
烹饪术语：汤煮久后水分蒸发、浓缩成精华（boil down 熬浓）。"熬到只剩核心"就是"归结为"。
\end{origin}

\vw{fine-tune}{/ˌfaɪn ˈtjuːn/}{vt.}

\sense{vt.}{对……微调}
\col{fine-tune a system/plan}
\ex{The engineers fine-tuned the engine before the race.}
\ex{We need to fine-tune the details of the agreement.}

\begin{origin}
20 世纪收音机调谐术语：fine（精细）+ tune（调音）——"精细地调准"。同源词：\textbf{tune}（曲调）、\textbf{tone}（音调）。tune 与 tone 都来自希腊语 \textit{tonos}（绷紧的弦）。
\end{origin}

\vw{initial}{/ɪˈnɪʃl/}{adj. \& n.}

\sense{adj.}{最初的，开始的}
\col{the initial stage/phase 最初阶段}
\ex{The initial results are encouraging.}
\col{initial investment 初始投资}

\sense{n.}{首字母}
\col{one's initials 姓名首字母}
\ex{His initials are J.K.}

\col{initially (adv.) 起初}
\ex{Initially, I didn't like the city.}

\begin{origin}
来自拉丁语 \textit{initialis}，源于 \textit{initium}（开端），词根 \textit{inire}（走入）：\textit{in-}（入）+ \textit{ire}（走）。同源词：\textbf{initiate}（发起）、\textbf{initiative}（主动权）。"走进来的第一步"就是开端。
\end{origin}

\vw{optimistic}{/ˌɒptɪˈmɪstɪk/}{adj.}

\sense{adj.}{乐观的，抱乐观看法的}
\col{be optimistic about sth}
\ex{She is optimistic about the future.}
\col{be optimistic that ...}
\ex{We are optimistic that the plan will work.}

\col{optimism (n.) 乐观主义；optimist (n.) 乐观主义者}

\begin{origin}
来自拉丁语 \textit{optimus}（最好的）——\textit{bonus}（好）的最高级。乐观主义者相信"最好的会到来"。同源词：\textbf{optimum}（最佳状态）。反义词 pessimistic（悲观的）源于拉丁人名 Pessimus（最坏的）。
\end{origin}

\vw{annoy}{/əˈnɔɪ/}{vt.}

\sense{vt.}{使恼怒，使生气}
\ex{His constant tapping annoyed everyone in the room.}
\col{be annoyed at/about sth 对……生气}
\ex{She was annoyed at the delay.}
\col{be annoyed with sb 生某人的气}
\ex{He was annoyed with himself for forgetting.}

\sense{vt.}{打扰，骚扰}
\ex{Stop annoying your sister!}

\begin{gram}[-ed 与 -ing 形容词]
\textit{annoyed}（感到恼怒的）与 \textit{annoying}（令人恼怒的）：-ed 形容人的感受，-ing 形容事物的性质——\textit{an annoying noise}（恼人的噪音）；\textit{I'm annoyed}（我很恼火）。同类：\textit{interested/interesting}、\textit{excited/exciting}、\textit{bored/boring}、\textit{frightened/frightening}。
\end{gram}

\begin{origin}
来自古法语 \textit{anoier}（使厌烦），源于拉丁语 \textit{inodiare}（使憎恶），词根 \textit{odium}（憎恨）。同源词：\textbf{odious}（可憎的）。"引起憎恶"就是使恼怒。
\end{origin}

\vw{weed}{/wiːd/}{n.}

\sense{n.}{杂草，野草}
\col{pull up weeds 拔草}
\ex{The garden is full of weeds.}
\col{weed (vt.) 除草（拓展）}
\ex{She spent the morning weeding the flower beds.}

\col{like weeds 大量地，像野草一样（疯长）}
\ex{New restaurants are sprouting up like weeds.}

\begin{origin}
源自古英语 \textit{wēod}（杂草），日耳曼语。野草蔓延极快，"像野草一样"便喻"大量涌现"——与 \textit{grow like mushrooms}（像蘑菇一样冒出来）同类比喻。
\end{origin}

\vw{state-of-the-art}{/ˌsteɪt əv ði ˈɑːt/}{adj.}

\sense{adj.}{使用最先进技术（或方法）的，最先进的}
\col{state-of-the-art technology/equipment}
\ex{The laboratory is equipped with state-of-the-art technology.}

\begin{origin}
"技艺的最高状态"——state（状态）+ of the + art（技艺，古义"手艺、工艺"，与"艺术"无涉）。20 世纪 60 年代工程界术语，指"当前工艺水平的极限"。
\end{origin}

\vw{globe}{/ɡləʊb/}{n.}

\sense{n.}{球体，球状物；地球仪}
\ex{The teacher pointed at the globe.}

\sense{n.}{地球，世界}
\col{across the globe 全世界}
\ex{People across the globe watched the event live.}
\col{global (adj.) 全球的}
\ex{global warming}

\begin{origin}
来自拉丁语 \textit{globus}（球）。同源词：\textbf{global}（全球的）、\textbf{globalization}（全球化）。"球"——地球本就是一个球体。
\end{origin}

\vw{expense}{/ɪkˈspens/}{n.}

\sense{n.}{费用；开销；开支}
\col{expenses 开支，开销}
\ex{The company covers all travel expenses.}
\col{a major expense 一笔大开销}

\col{at the expense of sth 在牺牲（或损害）……的情况下}
\ex{He pursued his career at the expense of his health.}
\col{at sb's expense 由某人出钱；拿某人开玩笑}
\ex{They had a good laugh at his expense.}
\col{spare no expense 不惜代价}
\ex{The couple spared no expense for the wedding.}

\begin{origin}
来自拉丁语 \textit{expendere}（称出、支付）：\textit{ex-}（出）+ \textit{pendere}（称、悬挂）——古时以"称"结算钱币。同源词：\textbf{expend}（花费）、\textbf{spend}（spend 源自 dispendere 的缩写）、\textbf{pension}（养老金）、\textbf{expensive}（昂贵的）。expense 与 spend 同根。
\end{origin}

\vw{pub}{/pʌb/}{n.}

\sense{n.}{酒吧，酒馆}
\ex{We met at the pub after work.}
\col{go to the pub 去酒吧}
\col{pub lunch 酒吧简餐}

\begin{origin}
pub 是 \textit{public house}（公共房屋）的缩略，17 世纪起指"向公众开放卖酒的房子"。英国文化中 pub 是社区生活中心，全英约有四万家。
\end{origin}

% ---------------- 选收的未标 * 词条 ----------------

\vw{analyse}{/ˈænəlaɪz/}{vt.}

\sense{vt.}{(AmE analyze) 分析}
\col{analyse data/results}
\ex{Scientists analysed the samples in the lab.}
\col{analysis (n.) 分析}
\ex{a detailed analysis of the situation}

\begin{origin}
来自希腊语 \textit{analyien}（解开）：\textit{ana-}（向上、分开）+ \textit{lyein}（松开）。同源词：\textbf{analysis}（分析）、\textbf{paralysis}（瘫痪：para- 旁 + lysis 松开 $\to$ 肢体"松掉"）。"拆开来看"就是分析——analysis 与 paralysis 共享 lyein 词根。
\end{origin}

\vw{greedy}{/ˈɡriːdi/}{adj.}

\sense{adj.}{贪心的，贪婪的}
\col{be greedy for sth}
\ex{He was greedy for power and money.}
\col{greed (n.) 贪婪}

\begin{origin}
源自古英语 \textit{grædig}，源于 \textit{græd}（贪婪），与"饥饿"的词根有关——同源德语 \textit{gierig}（贪婪的）。"像饿鬼一样想占有"就是贪婪。
\end{origin}

\vw{database}{/ˈdeɪtəbeɪs/}{n.}

\sense{n.}{数据库}
\col{a database of sth}
\ex{The police keep a database of fingerprints.}
\col{access a database 访问数据库}

\begin{origin}
data（数据，datum 的复数）+ base（基础）——20 世纪 60 年代计算机术语。"数据的地基"就是数据库。
\end{origin}

\vw{hacker}{/ˈhækə(r)/}{n.}

\sense{n.}{黑客}
\col{a computer hacker 电脑黑客}
\ex{Hackers broke into the company's system.}

\begin{origin}
hack（砍、劈）+ -er。20 世纪 50 年代美国麻省理工学院校园俚语中，hack 指"巧妙的恶作剧或程序"——hacker 本指"程序高手"，后转指"入侵系统的黑客"。与 Unit 4 的 \textbf{break into}（侵入系统）正是一对。
\end{origin}

\vw{theft}{/θeft/}{n.}

\sense{n.}{偷窃，偷窃罪}
\col{car theft 汽车盗窃}
\ex{He was arrested for theft.}

\begin{origin}
源自古英语 \textit{þēofth}，源于 \textit{þēof}（贼）。同源词：\textbf{thief}（贼）、\textbf{thieve}（偷窃）。"贼的行为"就是偷窃。
\end{origin}

\vw{drone}{/drəʊn/}{n.}

\sense{n.}{无人驾驶飞机}
\col{a military drone 军用无人机}
\ex{Drones are now used to deliver parcels.}

\begin{origin}
drone 本义"雄蜂"（古英语 \textit{drān}）——雄蜂只会嗡嗡飞，是拟声词。20 世纪无人靶机飞行时嗡嗡作响如雄蜂，遂得名 drone。从雄蜂到无人机，一词跨千年。
\end{origin}

\vw{illegal}{/ɪˈliːɡl/}{adj.}

\sense{adj.}{不合法的，非法的}
\col{illegal activities/immigration}
\ex{It is illegal to park here.}

\begin{origin}
\textit{il-}（不）+ legal（合法的）。legal 来自拉丁语 \textit{lex}（法律）。同源词：\textbf{legal}、\textbf{legislation}（立法）、\textbf{privilege}（privus 私 + lex 法 $\to$ 私法 $\to$ 特权）。反义词：legal、lawful。
\end{origin}

\vw{licence}{/ˈlaɪsns/}{n.}

\sense{n.}{(AmE license) 许可证，执照}
\col{a driving licence 驾照}
\ex{He lost his licence for speeding.}
\col{a licence to do sth 做……的许可}
\ex{The company has a licence to operate the service.}

\begin{origin}
来自拉丁语 \textit{licentia}（自由、许可），源于 \textit{licere}（被允许）。同源词：\textbf{illicit}（非法的：in- 不 + licere）。"被允许的凭证"就是执照——licence 与 illicit 同根。
\end{origin}

\vw{satellite}{/ˈsætəlaɪt/}{n.}

\sense{n.}{人造卫星；卫星}
\col{a communications satellite 通信卫星}
\ex{The satellite orbits the Earth.}
\col{satellite TV 卫星电视}

\begin{origin}
来自拉丁语 \textit{satelles}（随从、侍卫）。1610 年伽利略用 satellites 称呼绕木星运转的四颗卫星——"行星的随从"。此后 satellite 泛指卫星与人造卫星。
\end{origin}

\vw{gender}{/ˈdʒendə(r)/}{n.}

\sense{n.}{性别}
\col{gender equality 性别平等}
\ex{The survey covers people of all genders.}

\begin{origin}
来自拉丁语 \textit{genus}（种类、出生），词根可追溯到原始印欧语 *\textit{gen-}（生育）——与 Unit 3 的 \textbf{gene}（基因）同源。同源词：\textbf{generate}（产生）、\textbf{general}（总的：属全类的）、\textbf{genius}（天赋）。
\end{origin}

\vw{envy}{/ˈenvi/}{n. \& vt.}

\sense{n. \& vt.}{羡慕，忌妒}
\col{envy sb sth 羡慕某人某物}
\ex{I envy you your good luck.}
\col{be the envy of sb 是某人羡慕的对象}
\ex{Her success is the envy of her classmates.}
\col{green with envy 嫉妒得眼红}

\begin{origin}
来自拉丁语 \textit{invidia}（忌妒），源于 \textit{invidere}（不怀好意地看）：\textit{in-}（向）+ \textit{videre}（看）。同源词：\textbf{envious}（忌妒的）、\textbf{invidious}（招人反感的）。“斜着眼看别人”就是忌妒——envy 与 video 同根。
\end{origin}

\vw{clone}{/kləʊn/}{n. \& vt.}

\sense{n.}{好像一模一样的人，复制品；克隆动物（或植物）}
\ex{The new phone is a clone of the original design.}

\sense{vt.}{克隆，以无性繁殖技术复制}
\ex{Scientists cloned the sheep named Dolly.}

\begin{origin}
来自希腊语 \textit{klōn}（嫩枝）。园艺上剪下一根嫩枝插活，即得母株的完整复制品——1903 年英国植物学家将此操作命名为 clone。克隆的本义，原是"插枝"。
\end{origin}

\end{document}
